\chapter*{Executive summary}

\renewcommand{\arraystretch}{1.5}

\begin{longtable}{|m{\textwidth}|}
\hline    

\centering\textbf{\large Title of Project} \\
\hline    
\centering Development of a Hydrogel Extruder \\
\hline    

\centering\textbf{\large Objectives} \\
\hline    
To design and construct a syringe-based hydrogel extruder and 3-axis gantry capable of fabricating three-dimensional structures of moderate complexity. \\
\hline    


\centering\textbf{\large What is current practice and what are its limitations?} \\
\hline    
Hydrogel bioprinters are currently used for tissue engineering and biomedical research, but existing printers are limited by slow curing rates, poor print resolution, and restricted control over extrusion parameters. \\
\hline    

\centering\textbf{\large What is new in this project?} \\
\hline    
A modular hydrogel printer capable of controlled extrusion of multiple types of hydrogel bioinks and integrated curing, calibrated on a test gel to prove the concept. \\
\hline    

\centering\textbf{\large If the project is successful, how will it make a difference?} \\
\hline    
The printer provides researchers with a versatile research tool to explore and fine-tune printing parameters for different hydrogels. It will contribute new knowledge to hydrogel 3D printing methods and support biomedical applications such as tissue scaffolding and cell culture environment fabrication. \\
\hline    

\centering\textbf{\large What are the risks to the project being a success? Why is it expected to be successful?} \\
\hline    
Risks include budget and time constraints, and challenges in sourcing an affordable, accessible test hydrogel. Success is expected through a clear design methodology, calibration with surrogate materials, and careful documentation. \\
\hline    

\centering\textbf{\large What contributions have/will other students made/make?} \\
\hline    
This is the second iteration of the project. Current work has established a functional prototype and experimental calibration framework. Future students can extend the design with advanced curing methods and dual extrusion capabilities. \\
\hline 

\centering\textbf{\large Which aspects of the project will carry on after completion and why?} \\
\hline    
The use of the printer will continue in Stellenbosch University's Physiology Department, where it will be used as a tool for characterising different types of bioinks. \\
\hline    

\centering\textbf{\large What arrangements have been/will be made to expedite continuation?} \\
\hline    
Comprehensive documentation of mechanical, electronic, and firmware design has been produced, along with an experimental calibration methodology, enabling seamless future development and integration into research environments. \\
\hline    

\end{longtable}
